\documentclass[11pt]{article}
\usepackage[T1]{fontenc}
\usepackage[utf8]{inputenc}
\usepackage[a4paper,margin=1in]{geometry}
\usepackage{parskip}
\usepackage{microtype}
\usepackage[hidelinks]{hyperref}

\setlength{\emergencystretch}{2em}

\begin{document}
\thispagestyle{empty}

\noindent\textbf{Phuc Hao Do}\\
Faculty of Information Technology\\
Danang Architecture University, Danang, Vietnam\\
\href{mailto:haodp@dau.edu.vn}{haodp@dau.edu.vn}

\vspace{1.2em}
\noindent June 27, 2026

\vspace{1.2em}
\noindent\textbf{To:} The Guest Editors,\\
IEEE Transactions on Network Science and Engineering,\\
Special Issue on \emph{Tutorials and Surveys on AI and Quantum-enabled Learning for
Future Wireless Networks}\\
(Prof. Trung Q. Duong, Prof. Ruidong Li, Dr. Bhaskara Narottama, Prof. George K.
Karagiannidis)

\vspace{1.0em}
\noindent\textbf{Re:} Submission of the manuscript ``\emph{Quantum and Quantum-Inspired
Graph Neural Networks for Dynamic Wireless Network Topologies: A Tutorial and Survey with
a Reproducible LEO Case Study}''

\vspace{1.2em}
\noindent Dear Guest Editors,

I am pleased to submit the above manuscript for consideration in the Special Issue on
\emph{Tutorials and Surveys on AI and Quantum-enabled Learning for Future Wireless
Networks}. The paper is prepared as a combined \textbf{tutorial and survey}, the
contribution type the call solicits, and it speaks directly to two of the listed topics:
\emph{quantum graph neural networks} and \emph{data-driven quantum algorithms for dynamic
wireless network topologies}.

\textbf{What the paper offers.} The manuscript is organized around a single unifying idea:
classical graph neural networks, graph diffusion, and continuous-time quantum walks are
the same object, a propagation operator on the graph Laplacian, so the only thing that
changes between them is the spreading physics (diffusive for the heat kernel, ballistic
for the quantum walk). From this view the paper (i) surveys the classical, quantum,
quantum-inspired, and dynamic-graph branches and their use in wireless and
non-terrestrial networks, organized by operator rather than by author; (ii) teaches, with
pseudocode and a worked example, how to build a param-matched quantum-inspired GNN layer
that runs on classical hardware today; and (iii) demonstrates the framework on a fully
reproducible low Earth orbit (LEO) case study.

\textbf{The central finding is honest and, we believe, timely for this Special Issue.} On
a node-potential task over a dynamic LEO graph, global propagation beats local message
passing by a wide margin at every scale, and the quantum-inspired and classical operators
trade places as the constellation grows: the quantum-inspired walk loses at 66
satellites, ties at 264, and becomes the best operator at 1584 satellites, exactly the
diameter-driven trend its ballistic spreading predicts. We report this conditional,
scale-dependent quantum advantage with its caveats rather than as a blanket claim, and we
add a NISQ-noise study (PennyLane) showing how device noise erodes the mechanism. A
dedicated section on evaluation pitfalls, scale-coupled normalization and single-seed
reporting, distills a reusable checklist; we believe this methodological contribution is
useful to the community independently of the case study.

\textbf{Reproducibility.} Every figure and number in the paper is regenerated from seeded
runs by the companion code at
\href{https://github.com/haodpsut/qi-gnn-dynamic-wireless}{github.com/haodpsut/qi-gnn-dynamic-wireless}.
All experiments are param-matched and multi-seed.

\textbf{Originality and disclosure.} This manuscript is original, has not been published
before, and is not under consideration elsewhere. For full transparency: I have a
separate, accepted survey on AI and quantum technologies in non-terrestrial networks (in
the \emph{Journal of Communications and Information Networks}). The present paper is
distinct in scope and contribution; it is a focused tutorial on graph-learning operators
with an original reproducible case study, not a broad survey, and it does not reuse that
material. The LEO topology generator reuses my open-source simulator, while the
node-potential task and all results here are new.

I confirm there are no conflicts of interest with the Guest Editors that would prevent a
fair review, and I am happy to suggest or exclude reviewers if helpful.

Thank you for your time and consideration.

\vspace{1.4em}
\noindent Sincerely,

\vspace{0.4em}
\noindent\textbf{Phuc Hao Do}\\
Danang Architecture University

\end{document}
